%%%%%%%%%%%%%%%%%%%%%%%%%%%%%%%%%%%%%%%%%
% Important note:
% This template must be compiled with XeLaTeX, the below lines will ensure this
%!TEX TS-program = xelatex
%!TEX encoding = UTF-8 Unicode
%
%%%%%%%%%%%%%%%%%%%%%%%%%%%%%%%%%%%%%%%%%

%----------------------------------------------------------------------------------------
%	PACKAGES AND OTHER DOCUMENT CONFIGURATIONS
%----------------------------------------------------------------------------------------

\documentclass[11pt, letterpaper]{awesome-cv} % A4 paper size by default, use 'letterpaper' for US letter

\geometry{left=2cm, top=1.5cm, right=2cm, bottom=2cm, footskip=.5cm} % Configure page margins with geometry
 
\fontdir[fonts/] % Specify the location of the included fonts

% Color for highlights
%\colorlet{awesome}{awesome-skyblue} % Default colors include: awesome-emerald, awesome-skyblue, awesome-red, awesome-pink, awesome-orange, awesome-nephritis, awesome-concrete, awesome-darknight
%\definecolor{awesome}{HTML}{CA63A8} % Uncomment if you would like to specify your own color
\definecolor{awesome}{HTML}{0BA89A} % UCSB sea green

% Colors for text - uncomment and modify
%\definecolor{darktext}{HTML}{414141}
%\definecolor{text}{HTML}{414141}
%\definecolor{graytext}{HTML}{414141}
%\definecolor{lighttext}{HTML}{414141}

\renewcommand{\acvHeaderSocialSep}{\quad\textbar\quad} % If you would like to change the social information separator from a pipe (|) to something else

%----------------------------------------------------------------------------------------
%	PERSONAL INFORMATION
%	Comment any of the lines below if they are not required
%----------------------------------------------------------------------------------------

\name{Jason C.}{Wei}
\address{2320 Willet Way, Pleasanton, California, 94566, USA}
\mobile{(925) 352-7600}
\github{github.com/Jayweiii}
\homepage{jweii.com}
\email{jwei@ucsb.edu}
\linkedin{linkedin.com/in/jwei}
%\skype{skypeid}
%\stackoverflow{SOid}{SOname}
%\twitter{@twit}
%\linkedin{linkedin name}
%\reddit{reddit account}
%\xing{xing name}
%\extrainfo{test} % Other text you want to include on this line

\position{Mechanical Engineer{\enskip\cdotp\enskip}UCSB} % Your expertise/fields
\quote{``Creativity is all about connecting dots."} % A quote or statement

%----------------------------------------------------------------------------------------
%	RECIPIENT/POSITION/LETTER INFORMATION
%	All of the below lines must be filled out
%----------------------------------------------------------------------------------------

\recipient{Internship Recruitment Team}{Lucid Motors\\Newark, CA} % The company being applied to

\letterdate{\today} % The date on the letter, default is the date of compilation

\lettertitle{Thermal Systems Intern} % The title of the letter

\letteropening{Dear Lucid Motors,} % How the letter is opened

\letterclosing{Kind regards,} % How the letter is closed

% \letterenclosure[Attached]{Resume} % Any enclosures with the letter

\makecvfooter{\today}{Jason C. Wei~~~·~~~Cover Letter}{} % Specify the letter footer with 3 arguments: (<left>, <center>, <right>), leave any of these blank if they are not needed
  
%----------------------------------------------------------------------------------------

\begin{document}

\makecvheader % Print the header

\makelettertitle % Print the title

%----------------------------------------------------------------------------------------
%	LETTER CONTENT
%----------------------------------------------------------------------------------------

\begin{cvletter}

%------------------------------------------------

My 3D printing journey began during the beginning of the pandemic. When I heard about companies like Bauer, a producer of hockey equipment, shifting their production to medical face shields due to the PPE shortage, I knew I needed to contribute using the 3D printer I had just built. After working with local hospitals, I developed a custom 3D printed face shield design and started a large effort to print and donate them to ICUs. I taught over 20 families how to set up their first 3D printer, repaired any that broke, and sourced PET film material directly from China. After buying two more FDM printers, my home sounded like a mini manufacturing plant that ran 24/7. Since donating over 3000 face shields, I have been exploring ways to take advantage of 3D printing's fast turn-around time and limitless geometric possibilities to advance our world in other ways. As the 3D printer tool lead at the California NanoSystems Institute (CNSI) and the technical director at Gaucho Racing, the FSAE Electric team at UC Santa Barbara, I am delighted to submit this mechanical engineering intern application.

At CNSI, I manage, repair, and improve processes for our FDM, SLA, DLP, and PolyJet 3D printers. I support researchers throughout the entire design and manufacturing process to help their parts meet design specifications at the lowest cost. As a result, I have developed a thorough understanding of the different design and material limitations of our printers, and the ability to combine additive manufacturing, subtractive manufacturing (CNC \& manual machining), and chemical post-processes (such as acetone vapor smoothing) in a single part. Since 3D printers are also our most used tools, I have intricately repaired many of them, down to the PCB soldering level, in order to save repair costs and reduce downtime. 

Additionally, I have gotten to apply my 3D printing expertise in designing UCSB's Formula SAE Electric race car. In order to reduce weight while maximizing structural integrity, I have designed and printed many of our mounting solutions in carbon fiber filled nylon. Although 3D printed materials are indeed very optimized for low weight, high strength materials, my current ambitions are to combine 3D printed geometries, such as gyroids, with phase-change thermofluidics, to increase the efficiency of heat exchangers.

At Carbon3D, I will bring my extensive 3D printing, CAD, and machining experience to push the limits of 3D printing even further.

I look forward to discussing this internship opportunity and my qualifications with you in further detail. 

Thank you so much for your time and consideration.
%------------------------------------------------


\end{cvletter}

%----------------------------------------------------------------------------------------

\makeletterclosing % Print the signature and enclosures

\end{document}